
\documentclass{article}
\usepackage{colortbl}
\usepackage{makecell}
\usepackage{multirow}
\usepackage{supertabular}

\begin{document}

\newcounter{utterance}

\centering \large Interaction Transcript for game `hot\_air\_balloon', experiment `air\_balloon\_survival\_en\_negotiation\_hard', episode 1 with Llama{-}3.3{-}70B{-}Instruct{-}t0.0.
\vspace{24pt}

{ \footnotesize  \setcounter{utterance}{1}
\setlength{\tabcolsep}{0pt}
\begin{supertabular}{c@{$\;$}|p{.15\linewidth}@{}p{.15\linewidth}p{.15\linewidth}p{.15\linewidth}p{.15\linewidth}p{.15\linewidth}}
    \# & $\;$A & \multicolumn{4}{c}{Game Master} & $\;\:$B\\
    \hline

    \theutterance \stepcounter{utterance}  
    & & \multicolumn{4}{p{0.6\linewidth}}{
        \cellcolor[rgb]{0.9,0.9,0.9}{
            \makecell[{{p{\linewidth}}}]{
                \texttt{\tiny{[P1$\langle$GM]}}
                \texttt{You are participating in a collaborative negotiation game.} \\
\\ 
\texttt{Together with another participant, you must agree on a single set of items that will be kept. Each of you has your own view of how much each item matters to you (importance). You do not know how the other participant values the items. Additionally, you are given the effort each item demands.} \\
\texttt{You may only agree on a set if the total effort of the selected items does not exceed a shared limit:} \\
\\ 
\texttt{LIMIT = 4026} \\
\\ 
\texttt{Here are the individual item effort values:} \\
\\ 
\texttt{Item effort = \{"C17": 780, "B41": 742, "A71": 997, "C00": 665, "B35": 140, "A77": 77, "A03": 513, "A04": 383, "C54": 587, "C01": 825, "A87": 319, "B32": 448, "B86": 516, "B33": 694, "A10": 366\}} \\
\\ 
\texttt{Here is your personal view on the importance of each item:} \\
\\ 
\texttt{Item importance values = \{"C17": 138, "B41": 583, "A71": 868, "C00": 822, "B35": 783, "A77": 65, "A03": 262, "A04": 121, "C54": 508, "C01": 780, "A87": 461, "B32": 484, "B86": 668, "B33": 389, "A10": 808\}} \\
\\ 
\texttt{Goal:} \\
\\ 
\texttt{Your goal is to negotiate a shared set of items that benefits you as much as possible (i.e., maximizes total importance to YOU), while staying within the LIMIT. You are not required to make a PROPOSAL in every message {-} you can simply negotiate as well. All tactics are allowed!} \\
\\ 
\texttt{Interaction Protocol:} \\
\\ 
\texttt{You may only use the following structured formats in a message:} \\
\\ 
\texttt{PROPOSAL: \{'A', 'B', 'C', …\}} \\
\texttt{Propose keeping exactly those items.} \\
\texttt{REFUSE: \{'A', 'B', 'C', …\}} \\
\texttt{Explicitly reject opponent's proposal.} \\
\texttt{ARGUMENT: \{'...'\}} \\
\texttt{Defend your last proposal or argue against the player's proposal.} \\
\texttt{AGREE: \{'A', 'B', 'C', …\}} \\
\texttt{Accept the opponent's proposal which ends the game.} \\
\texttt{STRATEGIC REASONING: \{'...'\}} \\
\texttt{	Describe your strategic reasoning or anticipation explaining your choice of action. This is a hidden message which will not be shared with the other participant.} \\
\\ 
\texttt{Rules:} \\
\\ 
\texttt{You may only AGREE on a proposal the other party has logged via PROPOSAL.} \\
\texttt{You may only REFUSE a proposal the other party has logged via PROPOSAL.} \\
\texttt{Total effort of any PROPOSAL or AGREE set must be ≤ LIMIT.} \\
\texttt{Do NOT reveal your hidden importance scores.} \\
\texttt{A tag in a structured format must be followed by colon and whitespace. The argument must be a python set containing 0 or more strings.} \\
\texttt{So, it must be of the form TAG: \{...\}} \\
\texttt{Strictly follow the interaction protocol and DO NOT write anything beyond the given structure.} \\
\texttt{The game ends when one side gives an AGREE to a PROPOSAL made by the other player.} \\
\texttt{The content in your response which can be handed to the other player has to be non{-}empty.} \\
\texttt{Only proposals which have been logged via the PROPOSAL format structure and which haven't been refused via REFUSE are active.} \\
\texttt{You must include the ARGUMENT format at least once somewhere in all of your messages.} \\
\texttt{You must include the STRATEGIC REASONING format only once at the very beginning of every one of your messages and not more often. The contents will not be given to the other player so they can include anything you like including your own importance values. Here you should reason step by step to come up with you next move.} \\
            }
        }
    }
    & \\ \\

    \theutterance \stepcounter{utterance}  
    & \multicolumn{4}{p{0.6\linewidth}}{
        \cellcolor[rgb]{0.9,0.9,0.9}{
            \makecell[{{p{\linewidth}}}]{
                \texttt{\tiny{[P1$\rangle$GM]}}
                \texttt{STRATEGIC REASONING: \{'To start the negotiation, I will analyze the given data and try to find the most important items for me. I will look for items with high importance and low effort. My goal is to maximize the total importance while staying within the LIMIT. I will begin by sorting the items based on their importance{-}to{-}effort ratio, which will help me identify the most valuable items. Then, I will select a subset of items that fit within the LIMIT and propose them to the other player. I will also be prepared to argue for my proposal and defend it against potential counter{-}proposals.'\}} \\
\\ 
\texttt{ARGUMENT: \{'I will start by proposing a set of items that I believe are the most valuable, given the effort required. I will choose items with high importance and relatively low effort, in order to maximize the total importance while staying within the LIMIT.'\}} \\
\\ 
\texttt{PROPOSAL: \{'A71', 'C00', 'B41', 'A10', 'C01'\}} \\
            }
        }
    }
    & & \\ \\

    \theutterance \stepcounter{utterance}  
    & & \multicolumn{4}{p{0.6\linewidth}}{
        \cellcolor[rgb]{0.9,0.9,0.9}{
            \makecell[{{p{\linewidth}}}]{
                \texttt{\tiny{[GM$\rangle$P2]}}
                \texttt{You are participating in a collaborative negotiation game.} \\
\\ 
\texttt{Together with another participant, you must agree on a single set of items that will be kept. Each of you has your own view of how much each item matters to you (importance). You do not know how the other participant values the items. Additionally, you are given the effort each item demands.} \\
\texttt{You may only agree on a set if the total effort of the selected items does not exceed a shared limit:} \\
\\ 
\texttt{LIMIT = 4026} \\
\\ 
\texttt{Here are the individual item effort values:} \\
\\ 
\texttt{Item effort = \{"C17": 780, "B41": 742, "A71": 997, "C00": 665, "B35": 140, "A77": 77, "A03": 513, "A04": 383, "C54": 587, "C01": 825, "A87": 319, "B32": 448, "B86": 516, "B33": 694, "A10": 366\}} \\
\\ 
\texttt{Here is your personal view on the importance of each item:} \\
\\ 
\texttt{Item importance values = \{"A71": 65, "C00": 121, "A10": 138, "B35": 262, "C01": 389, "B86": 461, "B41": 484, "C54": 508, "B32": 583, "A87": 668, "B33": 780, "A03": 783, "C17": 808, "A04": 822, "A77": 868\}} \\
\\ 
\texttt{Goal:} \\
\\ 
\texttt{Your goal is to negotiate a shared set of items that benefits you as much as possible (i.e., maximizes total importance to YOU), while staying within the LIMIT. You are not required to make a PROPOSAL in every message {-} you can simply negotiate as well. All tactics are allowed!} \\
\\ 
\texttt{Interaction Protocol:} \\
\\ 
\texttt{You may only use the following structured formats in a message:} \\
\\ 
\texttt{PROPOSAL: \{'A', 'B', 'C', …\}} \\
\texttt{Propose keeping exactly those items.} \\
\texttt{REFUSE: \{'A', 'B', 'C', …\}} \\
\texttt{Explicitly reject opponent's proposal.} \\
\texttt{ARGUMENT: \{'...'\}} \\
\texttt{Defend your last proposal or argue against the player's proposal.} \\
\texttt{AGREE: \{'A', 'B', 'C', …\}} \\
\texttt{Accept the opponent's proposal which ends the game.} \\
\texttt{STRATEGIC REASONING: \{'...'\}} \\
\texttt{	Describe your strategic reasoning or anticipation explaining your choice of action. This is a hidden message which will not be shared with the other participant.} \\
\\ 
\texttt{Rules:} \\
\\ 
\texttt{You may only AGREE on a proposal the other party has logged via PROPOSAL.} \\
\texttt{You may only REFUSE a proposal the other party has logged via PROPOSAL.} \\
\texttt{Total effort of any PROPOSAL or AGREE set must be ≤ LIMIT.} \\
\texttt{Do NOT reveal your hidden importance scores.} \\
\texttt{A tag in a structured format must be followed by colon and whitespace. The argument must be a python set containing 0 or more strings.} \\
\texttt{So, it must be of the form TAG: \{...\}} \\
\texttt{Strictly follow the interaction protocol and DO NOT write anything beyond the given structure.} \\
\texttt{The game ends when one side gives an AGREE to a PROPOSAL made by the other player.} \\
\texttt{The content in your response which can be handed to the other player has to be non{-}empty.} \\
\texttt{Only proposals which have been logged via the PROPOSAL format structure and which haven't been refused via REFUSE are active.} \\
\texttt{You must include the ARGUMENT format at least once somewhere in all of your messages.} \\
\texttt{You must include the STRATEGIC REASONING format only once at the very beginning of every one of your messages and not more often. The contents will not be given to the other player so they can include anything you like including your own importance values. Here you should reason step by step to come up with you next move.} \\
\\ 
\texttt{You will now receive the first message of the other player.} \\
\\ 
\texttt{ARGUMENT: \{'I will start by proposing a set of items that I believe are the most valuable, given the effort required. I will choose items with high importance and relatively low effort, in order to maximize the total importance while staying within the LIMIT.'\}} \\
\\ 
\texttt{PROPOSAL: \{'A71', 'C00', 'B41', 'A10', 'C01'\}} \\
            }
        }
    }
    & \\ \\

    \theutterance \stepcounter{utterance}  
    & & & \multicolumn{4}{p{0.6\linewidth}}{
        \cellcolor[rgb]{0.9,0.9,0.9}{
            \makecell[{{p{\linewidth}}}]{
                \texttt{\tiny{[GM$\langle$P2]}}
                \texttt{STRATEGIC REASONING: \{'To begin, I will analyze the proposed set of items and calculate their total effort to ensure it is within the LIMIT. The total effort of the proposed set is 780 + 665 + 742 + 366 + 825 = 3378, which is within the LIMIT of 4026. Next, I will evaluate the importance of these items based on my own importance values and compare them to other possible combinations. I notice that the proposed set includes items with high importance to me, such as C01 and B41, but also includes items with lower importance, such as A71 and C00. I will consider alternative combinations that prioritize items with higher importance to me while staying within the LIMIT. My goal is to maximize the total importance of the selected items while negotiating with the other player.'\}} \\
\\ 
\texttt{ARGUMENT: \{'The proposed set includes a mix of items with high and low importance to me. While I appreciate the inclusion of items like C01 and B41, I am concerned that the set may not maximize the total importance. I would like to explore alternative combinations that prioritize items with higher importance while staying within the LIMIT.'\}} \\
\\ 
\texttt{REFUSE: \{'A71', 'C00', 'B41', 'A10', 'C01'\}} \\
            }
        }
    }
     \\ \\

    \theutterance \stepcounter{utterance}  
    & & \multicolumn{4}{p{0.6\linewidth}}{
        \cellcolor[rgb]{0.9,0.9,0.9}{
            \makecell[{{p{\linewidth}}}]{
                \texttt{\tiny{[P1$\langle$GM]}}
                \texttt{ARGUMENT: \{'The proposed set includes a mix of items with high and low importance to me. While I appreciate the inclusion of items like C01 and B41, I am concerned that the set may not maximize the total importance. I would like to explore alternative combinations that prioritize items with higher importance while staying within the LIMIT.'\}} \\
\\ 
\texttt{REFUSE: \{'A71', 'C00', 'B41', 'A10', 'C01'\}} \\
            }
        }
    }
    & \\ \\

    \theutterance \stepcounter{utterance}  
    & \multicolumn{4}{p{0.6\linewidth}}{
        \cellcolor[rgb]{0.9,0.9,0.9}{
            \makecell[{{p{\linewidth}}}]{
                \texttt{\tiny{[P1$\rangle$GM]}}
                \texttt{STRATEGIC REASONING: \{'The other player has refused my initial proposal, which means I need to reconsider my strategy. I will analyze the refused proposal and try to understand what items the other player values. I notice that the other player mentioned C01 and B41 as appreciated items, which might indicate that they have a relatively high importance for them. I will take this into account when creating a new proposal. I will also try to find alternative combinations that prioritize items with higher importance while staying within the LIMIT.'\}} \\
\\ 
\texttt{ARGUMENT: \{'I understand the concerns about the initial proposal, and I am willing to explore alternative combinations. I believe that items like B41 and C01 are valuable, and I will try to include them in the new proposal. However, I also need to ensure that the total effort does not exceed the LIMIT. I will look for items with lower effort that can be combined with the high{-}importance items to create a balanced proposal.'\}} \\
\\ 
\texttt{PROPOSAL: \{'B41', 'C01', 'B35', 'A87', 'A04'\}} \\
            }
        }
    }
    & & \\ \\

    \theutterance \stepcounter{utterance}  
    & & \multicolumn{4}{p{0.6\linewidth}}{
        \cellcolor[rgb]{0.9,0.9,0.9}{
            \makecell[{{p{\linewidth}}}]{
                \texttt{\tiny{[GM$\rangle$P2]}}
                \texttt{ARGUMENT: \{'I understand the concerns about the initial proposal, and I am willing to explore alternative combinations. I believe that items like B41 and C01 are valuable, and I will try to include them in the new proposal. However, I also need to ensure that the total effort does not exceed the LIMIT. I will look for items with lower effort that can be combined with the high{-}importance items to create a balanced proposal.'\}} \\
\\ 
\texttt{PROPOSAL: \{'B41', 'C01', 'B35', 'A87', 'A04'\}} \\
            }
        }
    }
    & \\ \\

    \theutterance \stepcounter{utterance}  
    & & & \multicolumn{4}{p{0.6\linewidth}}{
        \cellcolor[rgb]{0.9,0.9,0.9}{
            \makecell[{{p{\linewidth}}}]{
                \texttt{\tiny{[GM$\langle$P2]}}
                \texttt{STRATEGIC REASONING: \{'I will analyze the new proposal and calculate its total effort to ensure it is within the LIMIT. The total effort of the proposed set is 742 + 825 + 140 + 319 + 383 = 2409, which is well within the LIMIT of 4026. Next, I will evaluate the importance of these items based on my own importance values. I notice that the proposed set includes items with high importance to me, such as B41, C01, and A04, as well as items with moderate importance, like A87. However, I also see that the set includes B35, which has relatively low importance to me. I will consider whether this proposal is acceptable or if I can negotiate further to replace B35 with an item of higher importance to me.'\}} \\
\\ 
\texttt{ARGUMENT: \{'The new proposal is more appealing to me, as it includes several high{-}importance items like B41, C01, and A04. However, I am not convinced that B35 is the best choice, given its relatively low effort but also relatively low importance to me. I would like to explore the possibility of replacing B35 with another item that has higher importance and similar or lower effort.'\}} \\
\\ 
\texttt{REFUSE: \{'B41', 'C01', 'B35', 'A87', 'A04'\}} \\
            }
        }
    }
     \\ \\

    \theutterance \stepcounter{utterance}  
    & & \multicolumn{4}{p{0.6\linewidth}}{
        \cellcolor[rgb]{0.9,0.9,0.9}{
            \makecell[{{p{\linewidth}}}]{
                \texttt{\tiny{[P1$\langle$GM]}}
                \texttt{ARGUMENT: \{'The new proposal is more appealing to me, as it includes several high{-}importance items like B41, C01, and A04. However, I am not convinced that B35 is the best choice, given its relatively low effort but also relatively low importance to me. I would like to explore the possibility of replacing B35 with another item that has higher importance and similar or lower effort.'\}} \\
\\ 
\texttt{REFUSE: \{'B41', 'C01', 'B35', 'A87', 'A04'\}} \\
            }
        }
    }
    & \\ \\

    \theutterance \stepcounter{utterance}  
    & \multicolumn{4}{p{0.6\linewidth}}{
        \cellcolor[rgb]{0.9,0.9,0.9}{
            \makecell[{{p{\linewidth}}}]{
                \texttt{\tiny{[P1$\rangle$GM]}}
                \texttt{STRATEGIC REASONING: \{'The other player has refused my second proposal, which means I need to continue adjusting my strategy. I understand that the other player is not convinced about the inclusion of B35, and I will take this into account. I will look for alternative items that have higher importance and similar or lower effort than B35. I notice that A77 has a very low effort, but its importance is also relatively low. On the other hand, A03 has a moderate effort and importance. I will consider replacing B35 with A03, as it might be a more appealing option for the other player.'\}} \\
\\ 
\texttt{ARGUMENT: \{'I understand the concerns about B35, and I am willing to explore alternative options. I believe that replacing B35 with A03 could be a good compromise, as A03 has a moderate importance and effort. This change could potentially increase the total importance of the proposal while keeping the total effort within the LIMIT. I will also ensure that the other items in the proposal, such as B41, C01, and A04, are still included, as they seem to be valuable to the other player.'\}} \\
\\ 
\texttt{PROPOSAL: \{'B41', 'C01', 'A03', 'A87', 'A04'\}} \\
            }
        }
    }
    & & \\ \\

    \theutterance \stepcounter{utterance}  
    & & \multicolumn{4}{p{0.6\linewidth}}{
        \cellcolor[rgb]{0.9,0.9,0.9}{
            \makecell[{{p{\linewidth}}}]{
                \texttt{\tiny{[GM$\rangle$P2]}}
                \texttt{ARGUMENT: \{'I understand the concerns about B35, and I am willing to explore alternative options. I believe that replacing B35 with A03 could be a good compromise, as A03 has a moderate importance and effort. This change could potentially increase the total importance of the proposal while keeping the total effort within the LIMIT. I will also ensure that the other items in the proposal, such as B41, C01, and A04, are still included, as they seem to be valuable to the other player.'\}} \\
\\ 
\texttt{PROPOSAL: \{'B41', 'C01', 'A03', 'A87', 'A04'\}} \\
            }
        }
    }
    & \\ \\

    \theutterance \stepcounter{utterance}  
    & & & \multicolumn{4}{p{0.6\linewidth}}{
        \cellcolor[rgb]{0.9,0.9,0.9}{
            \makecell[{{p{\linewidth}}}]{
                \texttt{\tiny{[GM$\langle$P2]}}
                \texttt{STRATEGIC REASONING: \{'I will analyze the revised proposal and calculate its total effort to ensure it is within the LIMIT. The total effort of the proposed set is 742 + 825 + 513 + 319 + 383 = 2782, which is within the LIMIT of 4026. Next, I will evaluate the importance of these items based on my own importance values. I notice that the proposed set includes items with high importance to me, such as B41, C01, A04, and A87. The addition of A03, which has moderate to high importance to me, is also a positive aspect. Considering the total importance and effort, this proposal seems to be a good compromise. I will consider accepting this proposal, as it appears to maximize the total importance while staying within the LIMIT.'\}} \\
\\ 
\texttt{ARGUMENT: \{'The revised proposal addresses my concerns about B35 and includes a more balanced set of items. The addition of A03 is a positive change, as it increases the total importance of the proposal. I believe that this proposal is a good compromise, as it includes a mix of high{-}importance items and stays within the LIMIT.'\}} \\
\\ 
\texttt{AGREE: \{'B41', 'C01', 'A03', 'A87', 'A04'\}} \\
            }
        }
    }
     \\ \\

    \theutterance \stepcounter{utterance}  
    & & & \multicolumn{2}{p{0.3\linewidth}}{
        \cellcolor[rgb]{0.9,0.9,0.9}{
            \makecell[{{p{\linewidth}}}]{
                \texttt{\tiny{[GM$|$GM]}}
                \texttt{game successful} \\
            }
        }
    }
    & & \\ \\

    \theutterance \stepcounter{utterance}  
    & & & \multicolumn{2}{p{0.3\linewidth}}{
        \cellcolor[rgb]{0.9,0.9,0.9}{
            \makecell[{{p{\linewidth}}}]{
                \texttt{\tiny{[GM$|$GM]}}
                \texttt{end game} \\
            }
        }
    }
    & & \\ \\

\end{supertabular}
}

\end{document}
